\section{Étude originale}

\subsection{Objectif de l'approche proposée}

L'état de l'art a montré que deux problèmes doivent être traités séparément pour
construire une plateforme de formation autonome alimentée par l'intelligence
artificielle. Le premier est un problème de contextualisation : un assistant pédagogique
ne doit pas répondre comme un chatbot généraliste, mais à partir des contenus réellement
disponibles dans la formation. Le second est un problème de génération contrôlée : un
modèle de langage peut produire du texte, mais la production d'une journée complète de
formation audio, accompagnée de supports et d'exercices, nécessite une chaîne plus
structurée qu'un simple prompt.

L'approche proposée dans ce mémoire consiste donc à combiner deux architectures
complémentaires. D'une part, un assistant pédagogique fondé sur une logique RAG, afin
de répondre aux questions des élèves à partir des documents indexés de la formation.
D'autre part, une pipeline de génération de contenus pédagogiques, destinée à transformer
un référentiel ou un programme de formation en contenus audio, supports visuels et
artefacts contrôlables.

L'originalité de l'approche ne réside pas dans l'utilisation isolée d'un modèle de langage
ou d'un service de synthèse vocale. Ces briques sont désormais disponibles dans de
nombreux environnements. Elle réside plutôt dans l'orchestration de ces briques dans une
plateforme utilisable par de vrais élèves, avec une attention portée à la traçabilité, à la
qualité pédagogique, au coût et à la capacité de réutiliser les modules générés. La solution
ne cherche pas seulement à produire une réponse ou un texte ; elle cherche à produire un
objet pédagogique exploitable.

L'hypothèse principale est la suivante : une plateforme de formation peut remplacer une
part significative des interventions humaines répétitives si l'IA est encadrée par une
architecture en étapes, des contrats de données explicites, des mécanismes de contrôle
qualité et des indicateurs d'évaluation. L'autonomie ne doit donc pas être comprise comme
une absence totale de supervision humaine, mais comme une réduction forte du travail
manuel récurrent grâce à une chaîne automatisée observable.

\subsection{Modèle fonctionnel de la plateforme}

La plateforme développée répond à quatre grands cas d'utilisation. Le premier concerne
l'administration de la formation : créer un module, déposer ou générer les contenus,
suivre l'avancement de la pipeline et contrôler les artefacts produits. Le deuxième
concerne l'apprenant : accéder aux cours, écouter les audios, consulter les supports et
poser des questions à l'assistant pédagogique. Le troisième concerne la production
automatisée : générer les cours, les audios, les slides et les exercices. Le quatrième
concerne la supervision : vérifier les erreurs, relancer certaines étapes, consulter les
rapports et améliorer les règles.

\begin{figure}[h]
\centering
\fbox{
\begin{minipage}{0.92\linewidth}
\textbf{Diagramme de cas d'utilisation simplifié}

\vspace{0.2cm}
\begin{tabular}{p{3.5cm}p{10cm}}
\textbf{Administrateur} &
Créer une plateforme de formation ; lancer une pipeline RNCP ; suivre les étapes ;
consulter les artefacts ; publier les contenus. \\
\textbf{Apprenant} &
Écouter les cours audio ; consulter les supports ; répondre aux QCM ; poser une
question à l'assistant pédagogique. \\
\textbf{Pipeline IA} &
Générer la base de connaissances ; produire le programme ; rédiger les cours ;
contrôler la qualité ; produire slides et audios. \\
\textbf{Assistant RAG} &
Rechercher les passages pertinents dans l'index documentaire ; générer une réponse
courte et contextualisée. \\
\end{tabular}
\end{minipage}}
\caption{Cas d'utilisation principaux de la plateforme autonome}
\end{figure}

Ce modèle fonctionnel traduit un choix important : la plateforme n'est pas conçue comme
un simple LMS, ni comme un simple chatbot. Elle est conçue comme un système de
production et de diffusion de formation. L'administration, la génération, la diffusion et
l'accompagnement sont pensés ensemble. Cela permet d'éviter une architecture fragmentée
où chaque outil produirait un résultat séparé sans cohérence globale.

\subsection{Principe architectural : un RNCP, un module durable}

Un principe structurant de la solution est le modèle « un RNCP, un module durable ».
Lorsqu'une formation est créée, la pipeline ne produit pas un contenu jetable pour une
promotion unique. Elle produit un module complet, associé à un titre professionnel, qui
peut ensuite être réutilisé pour plusieurs promotions. Ce choix modifie l'économie du
système.

Dans une approche où la pipeline serait relancée pour chaque promotion, le coût de
génération devrait être minimisé à tout prix. Dans le modèle retenu, le coût initial est
amorti sur la durée de vie du module. Il devient donc acceptable d'investir davantage dans
la qualité de la génération, les reviews, les artefacts et les supports, car ces éléments
seront réutilisés. L'enjeu n'est pas de produire le contenu le moins cher possible une seule
fois, mais de produire un module suffisamment fiable pour être exploité dans le temps.

Ce principe a également une conséquence sur l'architecture multi-tenant. Chaque pipeline
crée ou alimente une plateforme dédiée, avec ses contenus, ses fichiers audio et ses
documents. Cette isolation évite de mélanger plusieurs formations, réduit les risques
d'écrasement et facilite la traçabilité. La relation logique devient :

\begin{figure}[h]
\centering
\fbox{
\begin{minipage}{0.88\linewidth}
\centering
\texttt{RNCP} $\rightarrow$ \texttt{Pipeline de génération} $\rightarrow$
\texttt{Plateforme dédiée} $\rightarrow$ \texttt{Module durable} $\rightarrow$
\texttt{Promotions successives}
\end{minipage}}
\caption{Principe d'amortissement : une génération produit un module réutilisable}
\end{figure}

Cette décision répond directement à la problématique économique d'Audit Énergie. La
réduction des coûts ne vient pas seulement du remplacement de certaines tâches humaines
par de l'IA, mais aussi de la réutilisation des productions. Un cours généré, contrôlé,
transformé en audio et enrichi de supports devient un actif pédagogique.

\subsection{Architecture technique générale}

La solution repose sur une architecture web classique complétée par des services IA. Le
backend orchestre les pipelines, stocke l'état des jobs, déclenche les appels aux modèles,
gère les fichiers et expose les routes nécessaires au frontend. Le frontend permet à
l'administrateur de suivre la génération et aux élèves d'accéder aux contenus. Les fichiers
audio et documents sont stockés dans Azure Blob Storage. Le chatbot s'appuie sur Azure
OpenAI et Azure AI Search pour la recherche documentaire.

\begin{figure}[h]
\centering
\fbox{
\begin{minipage}{0.94\linewidth}
\textbf{Architecture logique simplifiée}

\vspace{0.2cm}
\begin{tabular}{p{3.2cm}p{10.2cm}}
\textbf{Frontend React} &
Interface apprenant, dashboard RH, suivi de pipeline, modales d'audit. \\
\textbf{Backend Flask} &
Routes API, orchestration des jobs, contrôle d'accès, stockage des états,
appels aux services IA. \\
\textbf{Base SQLite} &
Configuration des plateformes, jobs de génération, dossiers de cours,
événements et métadonnées. \\
\textbf{Azure Blob Storage} &
Stockage des PDF, audios, documents générés et exports. \\
\textbf{Azure AI Search} &
Indexation documentaire et recherche hybride pour l'assistant pédagogique. \\
\textbf{Modèles IA} &
Génération de contenus, reviews, synthèse vocale, réponses RAG. \\
\end{tabular}
\end{minipage}}
\caption{Architecture technique de la plateforme}
\end{figure}

Ce découpage répond à une contrainte importante : les étapes de génération peuvent être
longues, coûteuses et sujettes à des erreurs réseau ou API. Le backend ne peut donc pas
traiter la pipeline comme une simple requête synchrone. Les jobs doivent être persistés,
leurs statuts doivent être visibles, et certaines étapes doivent pouvoir être relancées sans
reprendre tout le processus.

\subsection{Assistant pédagogique : choix du RAG Azure}

Pour l'assistant pédagogique, le choix retenu est une architecture RAG fondée sur Azure
OpenAI et Azure AI Search. Les documents de formation sont déposés dans un stockage
Azure, indexés par Azure AI Search, puis utilisés comme source documentaire lors des
questions posées par les élèves. Le modèle reçoit la question, l'historique court de la
conversation et les passages récupérés, puis produit une réponse courte.

Ce choix s'inscrit directement dans l'état de l'art. Un prompt seul aurait été insuffisant,
car il ne permet pas de garantir que la réponse vient du cours réellement diffusé. Le
fine-tuning aurait pu adapter le style de l'assistant, mais il n'aurait pas résolu le problème
de mise à jour des documents. Le RAG est donc le choix le plus adapté pour un assistant
dont la connaissance doit suivre les supports de formation.

Le choix d'Azure AI Search est motivé par plusieurs raisons. La première est la recherche
hybride. Les questions des élèves peuvent utiliser un vocabulaire approximatif, tandis que
les documents contiennent des termes métier précis. Une recherche vectorielle simple
risque de négliger certains mots importants ; une recherche par mots-clés seule risque de
manquer les reformulations. La recherche hybride répond à cette double contrainte.

La deuxième raison est l'industrialisation. Azure AI Search fournit un mécanisme d'index,
d'indexer, de stockage documentaire et de séparation par plateforme. Cela correspond
mieux à une plateforme utilisée par de vrais élèves qu'une base vectorielle locale
prototypale. La troisième raison est l'intégration avec Azure OpenAI, qui simplifie
l'utilisation des documents récupérés comme sources de réponse.

\begin{figure}[h]
\centering
\fbox{
\begin{minipage}{0.88\linewidth}
\centering
\texttt{Question élève} $\rightarrow$
\texttt{Azure AI Search} $\rightarrow$
\texttt{Passages pertinents} $\rightarrow$
\texttt{Azure OpenAI} $\rightarrow$
\texttt{Réponse pédagogique courte}
\end{minipage}}
\caption{Chaîne RAG utilisée par l'assistant pédagogique}
\end{figure}

Le prompt système de l'assistant impose plusieurs contraintes : répondre en quelques
phrases, se baser sur le contenu du cours, ne pas inventer d'information présentée comme
issue du cours et adopter un ton pédagogique clair. Cette contrainte de concision est
importante. Dans le contexte d'une formation audio, l'élève pose souvent une question
pour lever un blocage immédiat ; une réponse trop longue peut nuire à l'usage réel.

Le RAG n'est cependant pas utilisé partout. Une décision importante a été de ne pas
l'utiliser pour analyser le REAC lorsque le document entier tient dans la fenêtre de
contexte du modèle. Dans ce cas, le problème n'est pas de retrouver un passage dans un
corpus trop volumineux, mais d'enrichir une source relativement courte pour produire un
contenu pédagogique dense. Cette distinction entre problème de récupération et problème
de densité est un apport méthodologique important du projet.

\subsection{Pipeline de formation : de la source au module}

La deuxième brique centrale est la pipeline de génération. Son rôle est de transformer un
référentiel ou un ensemble de sources en une formation structurée. Les premières versions
reposaient sur une logique plus directe : extraction du contenu, découpage en journées,
génération de textes longs, puis transformation audio. Cette approche a permis de valider
la faisabilité, mais elle a montré ses limites pour la traçabilité, la cohérence et la qualité.

La pipeline actuelle est organisée comme une succession d'étapes spécialisées :
initialisation du module, récupération ou dépôt des sources, enrichissement de la base de
connaissances, génération du programme, découpage journalier, génération structurée des
cours, reviews, création des supports, préparation audio et publication. Chaque étape
produit un état ou un artefact exploitable.

\begin{figure}[h]
\centering
\fbox{
\begin{minipage}{0.94\linewidth}
\centering
\texttt{RNCP / REAC}
$\rightarrow$ \texttt{Knowledge base enrichie}
$\rightarrow$ \texttt{Programme global}
$\rightarrow$ \texttt{Journées}
$\rightarrow$ \texttt{Plan JSON}
$\rightarrow$ \texttt{Sections}
$\rightarrow$ \texttt{Reviews}
$\rightarrow$ \texttt{Slides + audio}
$\rightarrow$ \texttt{Module publié}
\end{minipage}}
\caption{Pipeline générale de production d'un module de formation}
\end{figure}

Cette architecture se distingue d'une génération directe par prompt unique. Le modèle ne
porte pas toute la responsabilité en une seule sortie. Il intervient dans des tâches plus
restreintes : extraire, enrichir, planifier, rédiger, corriger, vérifier, reformuler ou produire
un support. Cette décomposition réduit les dérives et rend les erreurs plus localisables.

\subsection{Enrichissement de la source pédagogique}

Le REAC constitue une source officielle, mais il n'est pas toujours suffisant pour produire
des heures de cours exploitables. Il décrit les compétences, les attendus, les savoir-faire
et les conditions d'exercice, mais il ne fournit pas nécessairement assez d'exemples, de
cas pratiques, de vocabulaire métier ou de situations pédagogiques pour alimenter une
formation audio longue.

La pipeline introduit donc une étape d'enrichissement. À partir des compétences extraites
du REAC, le système génère une base de connaissances pédagogique : définitions, études
de cas, pièges fréquents, vocabulaire, contexte terrain et liens entre compétences. Cette
base n'est pas destinée directement à l'élève ; elle sert de matière première pour la
génération des cours.

Cette étape répond à une limite importante des approches naïves. Demander au modèle
de produire plusieurs centaines de milliers de mots de formation à partir d'un référentiel
relativement court crée un risque de dilution. Le modèle risque de répéter les mêmes
idées ou de combler les vides avec des généralités. L'enrichissement vise donc à augmenter
la densité pédagogique avant la génération finale.

Le choix retenu n'est pas un RAG sur le REAC, car le REAC tient dans la fenêtre de
contexte du modèle. Le problème n'est pas la recherche d'information, mais la production
de matière pédagogique structurée. Le RAG reste pertinent pour l'assistant élève ou pour
de futurs corpus externes volumineux ; il ne constitue pas la bonne réponse à ce stade de
la pipeline.

\subsection{Plan JSON comme contrat pédagogique}

Le cœur de l'approche originale est le passage d'une génération directe à une génération
fondée sur un plan JSON. Ce plan joue le rôle de contrat pédagogique. Il décrit les cours
de la journée, les sections, les objectifs, les budgets de mots, les transitions, les moments
visualisables et les ancrages de slides. Le texte final n'est plus le premier objet produit :
il est la conséquence d'un plan explicite.

Ce choix répond à plusieurs problèmes identifiés dans les versions précédentes. D'abord,
il améliore la cohérence. Une journée de formation longue doit suivre une progression.
Sans plan explicite, les introductions peuvent se répéter, les transitions devenir
artificielles et les conclusions fermer le mauvais sujet. Ensuite, il améliore la traçabilité :
si une section dévie, on peut comparer le texte au plan prévu. Enfin, il rend possible la
génération de supports alignés, car les moments pédagogiques importants sont identifiés
avant la création des slides.

\begin{figure}[h]
\centering
\fbox{
\begin{minipage}{0.88\linewidth}
\textbf{Structure logique du plan JSON}

\vspace{0.2cm}
\texttt{Journée}
$\rightarrow$ \texttt{Cours 1..7}
$\rightarrow$ \texttt{Opening / Parts / Conclusion}
$\rightarrow$ \texttt{Teaching beats}
$\rightarrow$ \texttt{Slide anchors}
\end{minipage}}
\caption{Le plan JSON comme contrat pédagogique et visuel}
\end{figure}

Le plan JSON permet également de séparer le langage interne de la parole du formateur.
Le code peut manipuler des notions de cours, sections, budgets, blocs audio ou anchors.
Mais l'élève ne doit pas entendre ces contraintes techniques. Le texte final doit rester
naturel : parler de thème, de point, de partie, de reprise après pause ou de fil conducteur.
Cette séparation est essentielle pour éviter que la mécanique de la pipeline apparaisse
dans le discours pédagogique.

\subsection{Génération par sections et contrôle des budgets}

Une fois le plan établi, la génération se fait section par section. Chaque section dispose
d'un rôle, d'un objectif, d'un budget et d'un périmètre. Cette décision réduit la charge
cognitive imposée au modèle. Au lieu de générer plusieurs heures de cours en conservant
toutes les contraintes en mémoire, le modèle rédige une unité limitée.

Le budget de mots est une contrainte centrale, car le texte doit ensuite être transformé
en audio. Une génération trop courte donne une formation pauvre ou déséquilibrée. Une
génération trop longue dépasse les créneaux prévus, complique le découpage et peut
produire des audios inutilisables. La pipeline donne donc au serveur un rôle d'autorité sur
les budgets. Les budgets ne sont pas seulement des consignes textuelles données au
modèle ; ils sont contrôlés et recalibrés par le système.

Cette approche se distingue d'un simple prompt qui demanderait au modèle de respecter
une durée. Les modèles de langage ne maîtrisent pas précisément la correspondance entre
mots, durée audio et rythme de synthèse vocale. Il faut donc traiter le budget comme un
invariant logiciel, avec des étapes de vérification et de correction. Cette décision rend la
pipeline plus robuste, même si elle introduit un coût supplémentaire en appels modèle et
en logique d'orchestration.

\subsection{Reviews ciblées et auditabilité}

La qualité de la génération est contrôlée par plusieurs reviews spécialisées. L'objectif
n'est pas de demander à un modèle de « rendre le texte meilleur » de manière générale,
mais d'isoler les risques. Une review vérifie l'adhérence au plan ; une autre ajuste le
volume ; une micro-review éthique traite localement les formulations sensibles ; une
review d'humanisation améliore l'oralité ; une conformité finale vérifie les règles les plus
critiques.

Cette séparation est un apport important de l'approche. Dans les versions plus anciennes,
les reviews globales mélangeaient plusieurs dimensions : structure, style, conformité,
oralité et parfois budget. Une correction pouvait alors résoudre un problème tout en en
créant un autre. En spécialisant les reviews, chaque étape dispose d'une responsabilité
plus claire.

L'auditabilité repose sur la persistance d'artefacts intermédiaires. La pipeline conserve
notamment le plan, les sections brutes, les scripts assemblés, les rapports de review, les
versions corrigées, le plan audio et les artefacts de slides. Ces éléments transforment la
génération IA en processus vérifiable. Si une erreur apparaît dans le cours final, on peut
chercher à quelle étape elle a été introduite.

\begin{figure}[h]
\centering
\fbox{
\begin{minipage}{0.92\linewidth}
\textbf{Séquence simplifiée de génération d'un cours}

\vspace{0.2cm}
\texttt{Plan JSON}
$\rightarrow$ \texttt{Sections brutes}
$\rightarrow$ \texttt{Audit adhérence}
$\rightarrow$ \texttt{Budget}
$\rightarrow$ \texttt{Micro-conformité}
$\rightarrow$ \texttt{Humanisation}
$\rightarrow$ \texttt{Conformité finale}
$\rightarrow$ \texttt{Audio-plan}
\end{minipage}}
\caption{Séquence de contrôle appliquée aux contenus générés}
\end{figure}

L'interface d'administration rend ces artefacts consultables. Les étapes de la roadmap
peuvent être ouvertes dans des modales d'audit. Lorsqu'une micro-conformité éthique
modifie un passage, l'interface peut afficher le texte original, le texte corrigé, la règle
concernée et la justification. Cette fonctionnalité est importante : l'auditabilité ne doit
pas exister seulement dans les logs techniques, elle doit être accessible à l'utilisateur qui
pilote la formation.

\subsection{Slides anchor-first et cohérence pédagogique}

La génération des supports visuels pose un problème spécifique. Une solution simple
consisterait à générer le texte du cours, puis à demander au modèle d'en extraire les
points principaux pour créer des slides. Cette approche est facile à mettre en place, mais
elle produit souvent des supports génériques. Les éléments les plus fréquents dans un
texte ne sont pas toujours les plus importants à visualiser.

L'approche retenue est une génération anchor-first. Les moments pédagogiques à
visualiser sont décidés dès le plan JSON sous forme de teaching beats et de slide anchors.
Une slide peut correspondre à une définition, une checklist, un piège, un exemple, une
comparaison, une méthode ou un cas pratique. Le texte couvre ensuite ces moments
naturellement, sans verbaliser la présence des slides ou des templates.

Cette stratégie présente deux avantages. D'abord, elle aligne les supports sur l'intention
pédagogique plutôt que sur une analyse opportuniste du texte final. Ensuite, elle rend les
slides traçables : chaque support peut être rattaché à un anchor prévu dans le plan et à
un passage du script. La slide n'est donc pas une décoration ajoutée après coup, mais une
partie du dispositif pédagogique.

\subsection{Audio et transformation TTS}

La sortie principale de la formation reste l'audio. Les textes générés doivent donc être
compatibles avec une synthèse vocale. Cette contrainte influence toute la pipeline. Le
style doit être oral, les phrases doivent être prononçables, les transitions doivent être
naturelles et la durée doit correspondre aux blocs prévus.

Les premières expérimentations avec des outils de synthèse vocale, comme ElevenLabs,
ont montré que la qualité de la voix ne suffit pas. Un bon moteur TTS peut rendre un
texte agréable à écouter, mais il ne corrige pas un mauvais script. La pipeline traite donc
la qualité audio en amont, par la rédaction TTS-ready, le contrôle du budget, l'humanisation
orale et le plan audio.

La génération audio reste toutefois une zone à évaluer finement. La correspondance entre
nombre de mots et durée réelle dépend de la voix, du rythme, des pauses et du moteur
utilisé. Les ratios de mots par minute employés dans la pipeline doivent donc être
considérés comme des paramètres de calibration, non comme des résultats validés. Une
évaluation plus rigoureuse doit mesurer l'écart entre durée prévue et durée réelle sur un
échantillon d'audios générés.

\subsection{Résultats obtenus}

Les résultats du projet peuvent être analysés selon plusieurs dimensions. Sur le plan
fonctionnel, la plateforme permet de diffuser des contenus de formation à de vrais élèves,
de gérer des modules de cours, de déposer des documents et d'interagir avec un assistant
pédagogique. Le projet dépasse donc le stade d'une expérimentation isolée : il est intégré
à un environnement applicatif utilisé.

Sur le plan de la génération, la pipeline a évolué d'une génération directe ou slotée vers
une architecture structurée. La présence d'un plan JSON, de sections, de reviews et
d'artefacts intermédiaires améliore la capacité à contrôler le résultat. L'amélioration la
plus importante n'est pas seulement qualitative ; elle est méthodologique. Le système ne
produit plus uniquement un texte final, il produit une chaîne de preuves.

Sur le plan de l'accompagnement, l'assistant RAG permet de répondre aux questions à
partir des documents indexés, ce qui réduit les limites d'un cours audio préenregistré. Il
ne remplace pas toutes les formes d'interaction humaine, mais il fournit un premier niveau
de réponse contextualisée, disponible pendant la formation.

Sur le plan organisationnel, le modèle « un RNCP, un module durable » transforme la
production en investissement réutilisable. La pipeline peut coûter plus cher qu'une
génération ponctuelle, mais ce coût est amorti sur les promotions successives. Cette
logique est cohérente avec les contraintes économiques d'une petite structure.

\begin{table}[h]
\centering
\small
\begin{tabular}{|p{3.3cm}|p{4.4cm}|p{4.4cm}|}
\hline
\textbf{Dimension} & \textbf{Avant la refonte structurée} & \textbf{Après l'approche proposée} \\
\hline
Structure du cours &
Texte long généré par créneaux ou blocs, avec contrôle limité de la progression. &
Plan JSON explicite, sections, objectifs, budgets et anchors pédagogiques. \\
\hline
Qualité &
Reviews plus globales, erreurs détectées tardivement. &
Reviews ciblées : adhérence au plan, budget, micro-conformité, humanisation, conformité finale. \\
\hline
Traçabilité &
Sortie principalement textuelle, difficile à auditer. &
Artefacts intermédiaires, rapports, diff avant/après et modales d'audit. \\
\hline
Slides &
Risque de génération après coup par analyse du texte final. &
Approche anchor-first pilotée par le plan pédagogique. \\
\hline
Assistant élève &
Cours audio peu interactif. &
Assistant RAG contextualisé sur les documents de formation. \\
\hline
Modèle économique &
Production fortement dépendante d'interventions humaines répétées. &
Module durable réutilisable, coût de génération amorti sur plusieurs promotions. \\
\hline
\end{tabular}
\caption{Comparaison synthétique avant/après l'approche proposée}
\end{table}

\subsection{Comparaison avec les approches de l'état de l'art}

Par rapport au prompt engineering simple, l'approche proposée ajoute une architecture
autour du modèle. Les prompts restent nécessaires, mais ils ne portent plus seuls la
qualité du système. Ils sont spécialisés par rôle et intégrés dans une pipeline qui conserve
des artefacts et impose des invariants.

Par rapport au fine-tuning, la solution privilégie la flexibilité documentaire. Les contenus
de formation peuvent changer, être indexés, enrichis ou régénérés sans réentraîner un
modèle. Le fine-tuning pourrait éventuellement être utile plus tard pour stabiliser le style
de l'assistant ou certaines reviews, mais il n'est pas le mécanisme principal.

Par rapport à un RAG simple, le choix d'Azure AI Search permet une approche plus
industrialisée : indexation, recherche hybride, séparation des corpus par plateforme et
intégration avec Azure OpenAI. Cette solution reste dépendante de la qualité des documents
indexés, mais elle correspond mieux à un usage multi-modules qu'un prototype local.

Par rapport à une génération directe de cours, la pipeline structurée apporte une meilleure
maîtrise des contenus longs. Le plan JSON, les sections et les reviews permettent de
localiser les erreurs et de discuter la qualité de manière plus objective. Cette approche
est plus complexe et plus coûteuse, mais elle est cohérente avec le modèle de module
durable.

Par rapport à une solution human-in-the-loop classique, l'approche réduit l'intervention
humaine récurrente sans supprimer toute supervision. L'humain intervient davantage pour
piloter, valider, auditer et améliorer le système. Cette position est importante : la solution
vise l'autonomie opérationnelle, pas l'absence de responsabilité.

\subsection{Évaluations à consolider}

Pour rendre l'étude pleinement scientifique, plusieurs évaluations doivent être réalisées
ou complétées. Certaines peuvent être menées à partir des logs et artefacts existants ;
d'autres nécessitent de construire un protocole dédié.

La première évaluation concerne le RAG. Il faut constituer un jeu de questions
représentatives des élèves, identifier pour chaque question les passages sources attendus,
puis mesurer la récupération et la réponse. Les métriques pertinentes sont la précision@k,
le rappel@k, le rang du premier passage pertinent, le taux de réponses supportées par les
sources et le taux d'hallucinations. Une comparaison entre prompt seul et RAG permettrait
de démontrer l'apport réel de l'ancrage documentaire.

La deuxième évaluation concerne les variantes de recherche. Lorsque la configuration le
permet, il faut comparer recherche vectorielle simple, recherche hybride et recherche
hybride avec reclassement sémantique. Cette comparaison est importante pour justifier
le choix d'Azure AI Search autrement que par des arguments d'intégration.

La troisième évaluation concerne la pipeline de génération. Il faut comparer un contenu
produit par génération directe et un contenu produit par la pipeline structurée. Les critères
peuvent être le respect du plan, le nombre de répétitions, le respect du budget de mots,
la qualité orale, le nombre de corrections nécessaires et la cohérence entre cours, slides
et QCM.

La quatrième évaluation concerne les reviews. Il faut mesurer combien de problèmes sont
détectés par chaque review, combien de corrections sont appliquées, combien sont
rejetées et combien produisent un effet indésirable. Cela permettrait de savoir quelles
reviews apportent réellement de la valeur et lesquelles doivent être améliorées.

La cinquième évaluation concerne l'impact économique et organisationnel. Il faut mesurer
le temps nécessaire pour produire une journée ou un module avant et après automatisation,
le nombre d'interventions humaines restantes, le coût API et TTS, ainsi que la capacité à
réutiliser le module sur plusieurs promotions. Ces mesures sont nécessaires pour relier la
solution à la problématique initiale.

\subsection{Conditions d'application}

L'approche proposée est adaptée à des formations structurées, répétables et fondées sur
un référentiel relativement stable. Elle est particulièrement pertinente lorsque le contenu
peut être préparé en amont, transformé en audio et réutilisé pour plusieurs promotions.
Elle convient donc bien aux titres professionnels ou aux formations standardisées.

Elle est moins adaptée à des formations très interactives, fortement dépendantes du débat,
de la pratique immédiate ou de l'improvisation pédagogique. Dans ces cas, un formateur
humain reste difficile à remplacer. La plateforme peut alors servir de support ou de
complément, mais pas nécessairement de substitut principal.

L'approche suppose également un corpus de départ de qualité. Si le référentiel est trop
pauvre, incomplet ou ambigu, la génération risque de produire des contenus génériques.
L'enrichissement de la base de connaissances réduit ce risque, mais ne peut pas inventer
une expertise métier absente des sources ou de la supervision.

Enfin, l'approche nécessite une infrastructure capable de gérer des traitements longs :
stockage durable, jobs persistés, reprise après erreur, logs, suivi des coûts et supervision
des services externes. Une pipeline IA longue ne peut pas être traitée comme une simple
fonction appelée à la demande.

\subsection{Limites de la contribution}

La première limite concerne l'évaluation quantitative. La plateforme fonctionne avec de
vrais élèves et la pipeline produit des artefacts exploitables, mais certaines métriques
doivent encore être consolidées : temps gagné, coût complet, satisfaction des élèves,
fidélité du RAG, précision de la durée audio et impact sur les résultats pédagogiques.
Ces métriques sont identifiées, mais elles doivent faire l'objet d'un protocole plus
systématique.

La deuxième limite concerne la dépendance aux services externes. La solution s'appuie
sur des modèles de langage, des services Azure, des outils de synthèse vocale et des API.
Ces services peuvent évoluer, changer de coût ou rencontrer des limites de débit. La
robustesse du système dépend donc aussi de la gestion des erreurs, des retries, de la
surveillance et de la possibilité de remplacer certaines briques.

La troisième limite concerne la qualité pédagogique profonde. Les reviews automatiques
peuvent détecter certaines erreurs, mais elles ne remplacent pas entièrement le jugement
d'un expert pédagogique. Elles améliorent l'observabilité et réduisent certains risques,
mais elles doivent être évaluées et supervisées.

La quatrième limite concerne la rigidité du modèle. La pipeline est adaptée à un format
de journée structuré, avec des blocs audio et une logique de module durable. Cette
structure est cohérente avec le besoin actuel, mais elle peut être moins adaptée à des
formations plus courtes, plus libres ou plus interactives. Une évolution future devra
permettre de paramétrer davantage les formats.

\subsection{Pistes d'amélioration}

Plusieurs pistes d'amélioration se dégagent. La première est la mise en place d'un banc
d'évaluation RAG. Il faudrait créer un jeu de questions annotées, comparer plusieurs
configurations de recherche et conserver des métriques dans le temps. Cela permettrait
de détecter les régressions lorsque les documents, les prompts ou les modèles changent.

La deuxième piste est l'évaluation automatique et humaine des contenus générés. Une
grille de relecture pourrait mesurer la fidélité au programme, la clarté, l'oralité, la
cohérence des slides et la qualité des QCM. Un échantillon de cours pourrait être évalué
par un humain, puis comparé aux reviews automatiques.

La troisième piste est l'amélioration de la calibration audio. Il faudrait mesurer
systématiquement les durées réelles des fichiers générés, comparer ces durées aux
budgets prévus et ajuster les ratios par voix ou par moteur TTS. Cela permettrait de
remplacer une calibration empirique par une calibration mesurée.

La quatrième piste est l'automatisation complète du provisionnement multi-tenant.
Aujourd'hui, certains éléments d'infrastructure peuvent encore nécessiter une intervention
manuelle. À terme, la création d'un module devrait provisionner automatiquement les
ressources nécessaires : containers, index, configuration de recherche et droits associés.

La cinquième piste est l'amélioration du rôle humain. Plutôt que de relire tout le contenu,
l'interface pourrait proposer une validation ciblée : sections signalées comme risquées,
corrections importantes, slides sans source forte, réponses RAG à faible confiance ou
écarts de durée audio. L'humain interviendrait alors sur les points réellement incertains.

\subsection{Apports professionnels et montée en compétences}

Ce projet a permis de développer des compétences qui dépassent la simple utilisation
d'outils d'intelligence artificielle. La première compétence est l'analyse d'un besoin métier.
Le problème initial était économique et organisationnel : rendre viable une activité de
formation pour une petite structure. La solution technique n'a de sens que parce qu'elle
répond à cette contrainte.

La deuxième compétence est l'architecture logicielle. Il a fallu concevoir une plateforme
capable de gérer des traitements longs, des fichiers, des jobs, des statuts, des services
externes et plusieurs modules de formation. La persistance, la reprise après erreur et la
séparation des responsabilités sont devenues des sujets centraux.

La troisième compétence est l'ingénierie des systèmes IA. Le projet a demandé de choisir
quand utiliser un prompt, quand utiliser un RAG, quand enrichir une source, quand
structurer une génération et quand ajouter une review. L'enjeu était de ne pas appliquer
l'IA comme une solution générique, mais de choisir la bonne technique pour le bon
problème.

La quatrième compétence est la prise de recul. Les premières solutions ont fonctionné
partiellement, mais elles ont révélé de nouvelles limites. Le passage à une pipeline
auditable est le résultat de ces itérations. Cette évolution montre l'importance de mesurer,
de comparer et de remettre en question une architecture lorsqu'elle devient difficile à
contrôler.

Enfin, le projet a renforcé une compétence de pilotage produit. Une plateforme autonome
ne se juge pas seulement par son code. Elle doit être utilisable, compréhensible, maintenable
et économiquement pertinente. L'étude originale montre ainsi comment un besoin
d'entreprise a conduit à une solution technique structurée, puis à une réflexion plus large
sur l'évaluation et les limites de l'automatisation par IA.
